% ==================== SECTION RESISTANCES ====================

\element{resistances}{
    \begin{question}{resist-1}
        En supposant que toutes les variables sont correctement déclarées, en considérant le programme ci-dessous : compléter le tableau donnant les résistances équivalentes calculées par le programme.

        \begin{minipage}{.55\textwidth}
            \lstinputlisting{sources/codes/resistances.c}
        \end{minipage}\hfill
        \begin{minipage}{.4\textwidth}
            {
                \Large
                \begin{tabular}{|l|c|}
                    \hline
                    \textbf{Calcul} & \textbf{Résultat ($\Omega$)} \\\hline
                    Série R1+R2 &  \\\hline
                    Parallèle R1//R2 &  \\\hline
                    Série R1+R2+R3 &  \\\hline
                    Parallèle R1//R2//R3 &  \\\hline
                \end{tabular}
            }
        \end{minipage}

        \evaluationProf
    \end{question}

    \begin{question}{resist-2}
        Analyser le code de ce programme en apportant une critique de son organisation et de sa structure.

        \textit{Indices : répétitions, modularité, réutilisabilité...}

        \evaluationProfGrille{2}{4}
    \end{question}

    \begin{question}{resist-3}
        Proposer une refactorisation du code en créant au minimum deux fonctions appropriées pour améliorer la structure du programme.
jkljkljklj
        \textbf{Écrire uniquement les prototypes et les définitions des fonctions (pas besoin de réécrire le main complet).}

        \evaluationProfGrille{2}{1}
    \end{question}
}
